\subsection{Simulation Environment}
\label{subsec:env}
% Facts in this subsection are taken from
% simulator/third_party/bluesky-gym/bluesky_zoo/sector_cr/sector_cr.py,
% marl/interfaces/bluesky_interface.py and marl/interfaces/observation.py.
We use the open-source BlueSky air traffic simulator~\cite{bluesky} through the
multi-agent \texttt{sector\_cr} (sector conflict resolution) environment of a
PettingZoo~\cite{pettingzoo} fork of BlueSky-Gym~\cite{bluesky-gym}.

\textbf{Scenario.} At the start of an episode a polygonal sector is generated at random:
vertices are sampled on a circle whose area is 3750~NM$^2$ until the polygon they span has
an area of at least 2400~NM$^2$, so the sector area lies between 2400 and 3750~NM$^2$. The
$n$ aircraft are placed at uniformly random positions inside the sector, subject to a minimum
initial separation of 6~NM between every pair (rejection sampling; a new sector is drawn if
the fleet cannot be placed). The original environment has no such constraint and can start an
episode with pairs already closer than the 5~NM separation minimum; with the constraint no
episode starts in loss of separation, so every loss of separation that is counted occurred
under the policy being evaluated. Each aircraft is
assigned its own exit waypoint: a point is sampled on the sector boundary and moved outward
along the ray from the sector center to ten times its distance, and the aircraft is
initialized with its heading toward that waypoint. All aircraft use BlueSky's point-mass
performance model of an Airbus A320, are created at the same flight level (FL350) and with
the same initial speed, and remain at that level, so conflicts are resolved in the
horizontal plane. The fleet size $n$ is fixed within an episode.

\textbf{Airspace scaling.} We use two rules for the sector size. In the \emph{fixed} airspace
the sector area range is the original one for every $n$, so the traffic density grows in
proportion to $n$. In the \emph{constant-density} airspace the area range is multiplied by
$n/n_{\text{ref}}$ with $n_{\text{ref}}=9$, the largest training fleet size, so the density
equals that of the fixed airspace at $n=9$ for every $n$. Training uses the fixed airspace;
the zero-shot evaluation is run in both (Section~\ref{subsec:protocol}). The two rules
separate two effects that the fixed airspace confounds: more vehicles, and more vehicles per
unit area.
% TODO(authors): verify altitude and speed units. The source defines ALTITUDE = 350 ("In FL")
% and AC_SPD = 150 and passes both directly to bs.traf.cre(); whether that call interprets
% them as FL/kt or m/(m/s) should be checked before the paper states an initial speed value.
% TODO(authors): deferred per decision 2026-11-25 (GAP-07/GAP-09/GAP-12) -- no static obstacles
% and no non-cooperative intruders are present; all aircraft are cooperating agents.

\textbf{Actions.} The environment applies an action as a change of the current heading and
of the current speed and sends the result to the simulator as BlueSky \texttt{HDG} and
\texttt{SPD} commands. The learned policy uses 15 discrete actions, the product of five
heading changes, $\{-22.5^\circ,\,-11.25^\circ,\,0^\circ,\,+11.25^\circ,\,+22.5^\circ\}$, and
three speed changes, $\{-\Delta v,\,0,\,+\Delta v\}$, where $\Delta v$ is the environment's
speed increment. One control step applies the commands and advances the simulation by five
steps of 1~s, that is, 5~s of simulated time.
% TODO(authors): state the value and unit of \Delta v. The source sets D_VELOCITY = 20/3 with
% the comment "kts", but adds it to bs.traf.cas before multiplying by the m/s-to-kt factor.
% TODO(authors): deferred per decision 2026-11-25 (GAP-07/GAP-09/GAP-12) -- continuous actions
% and thrust/pitch/roll/yaw-rate (six-degree-of-freedom) control are not used.

\textbf{Observations and state.} The quantities of Section~\ref{subsec:entities} are computed
from the simulator's aircraft positions and ground-velocity components. Relative positions
use a local flat-earth frame ($x$ east, $y$ north). Positions and distances are divided by
50~km and velocities by 150~m/s. The own features $(\cos\chi_i,\sin\chi_i,v_i)$ are taken
from the environment's own observation, in which the airspeed is centered and scaled. The
environment's native observation, which lists only the three nearest aircraft, is not used.
With $r_{\text{detect}}=20$~NM the sensor radius is four times the separation minimum.

\textbf{Reward.} The reward of aircraft $i$ at a control step is the sum of a drift penalty
and an intrusion penalty,
\begin{equation}
\label{eq:env-reward}
r_i = -0.1\,|\chi_i| \;-\; \big|\{\, j\neq i \;:\; d_{ij} < 5\ \text{NM} \,\}\big| ,
\end{equation}
with the drift angle $\chi_i$ in radians, wrapped to $[-\pi,\pi]$. The team reward is the
mean of $r_i$ over the aircraft (Eq.~\eqref{eq:team-reward}).

\textbf{Episode end.} The environment has no terminal state. It truncates an episode when
all aircraft are outside the sector. Because the environment has no time limit of its own,
we add a limit of 150 control steps (750~s of simulated time). Since both conditions are
truncations, the target of Eq.~\eqref{eq:target} bootstraps on every transition ($e=0$).

\textbf{Seeding.} The scenario generator draws from the global NumPy random number
generator, which we seed explicitly in each environment process; the seed argument of the
environment's reset function has no effect.
